%!TEX encoding = UTF8
%!TEX root = 0-notes.tex

\unchapter{Fonctions}
\fancyhead[C]{\textbf{Chapitre \thechapter~— Fonctions}}
\addcontentsline{toc}{section}{Méthodologie}


\dfn{fonction, variable, constante}{
	On dit qu'une quantité $y\in\R$ s'exprime \emphindex{en fonction d}'une quantité $x\in\R$ lorsque, à chaque nombre $x$ possible, on peut associer \underline{une unique} valeur $y$.
	
	On note alors $y = f(x)$, lu  « $y$ égal $f$ de $x$ », où $x$ est la \emphindex{variable}, et $f$ la \emphindex{fonction}.
	
	Si une quantité ne dépend pas de $x$, on l'appelle \emphindex{constante}, ou \emphindex{indépendante} de $x$.
	Cependant, une constante peut toujours être vue comme une fonction de $x$ !
	
%	\begin{center}
%	\includegraphics[page=19]{2-figures.pdf}
%	\end{center}
}{}

\dfn{antécédent, image}{
	Considérons une fonction $f$ et deux nombres réels $x, y\in\R$ vérifiants
		\[ y = f(x). \]
	On lit \og $y$ égal $f$ de $x$ \fg, et on dit alors que
		\begin{enumerate}
			\item
			$y$ est l'\emphindex{image} de $x$ par $f$ ; et
			\item
			$x$ est \underline{un} \emphindex{antécédent} de $y$ par $f$.
		\end{enumerate}
}{}


\dfn{forme algébrique}{
	On appelle 
		\[ f(x) = 3 x+1 \]
	la \emphindex{forme algébrique} de $f$.
	On lit dans ce cas \og $f$ de $x$ égal trois $x$ plus 1 \fg.
	
	C'est une forme entièrement générale qui permet de déduire l'image $f(x)$ à partir de n'importe quel réel $x\in\D$ du domaine de $f$.
}{}


\dfn{domaine de définition}{
	Le domaine de définition d'une fonction $f$ est le plus grand ensemble d'antécédents $x$ sur lequel l'image $f(x)$ est bien définie.
}{}


\dfn{intervalle}{
	Soient $a, b \in \R$.
	L'\emphindex{intervalle} $[a;b]$ est l'ensemble des nombres réels compris entre $a$ et $b$ inclus.
		\[ x \in [a ; b] \qquad \iff \qquad x \in \R, a \leq x \leq b. \]
	Il faut donc que $a \leq b$ pour que l'intervalle soit non vide.
	
	Lorsqu'un crochet de l'intervalle regarde vers l'extérieur, la borne correspondante est exclue.
	
	Les symboles $a$ et $b$ peuvent prendre la notation $\minfty$ ou $\pinfty$ respectivement.
	L'infini n'est jamais inclu car ce n'est pas un nombre réel.
}{}

%\notations{
%	On définit les signes suivants correspondant à des inégalités \emphindex{strictes} et \emphindex{larges}.
%	\begin{multicols}{2}
%	\begin{enumerate}[leftmargin=50pt]
%		\item[$<$ :] strictement inférieur à
%		\item[$\leq$ :] inférieur ou égal à
%		\item[$>$ :] strictement supérieur à
%		\item[$\geq$ :] supérieur ou égal à
%	\end{enumerate}
%	\end{multicols}
%	On attire l'attention sur le fait que le « ou » est inclusif : $x \leq x$ est toujours vrai. 
%}

%\notations{
%	On peut noter deux inégalités sur la même lignes dès qu'elles vont dans le même sens.
%	Par exemple, les deux inégalités ci-dessus peuvent être condensées en une :
%		\[ x \in \R, \text{ et } a \leq x \leq b, \]
%	ou encore
%		\[ x \in \R, \text{ et } b \geq x \geq a. \]
%}


\dfn{courbe représentative}{
	Considérons $f$ une fonction sur un domaine $\D$.
	
	La \emphindex{courbe représentative} de $f$, notée $\Ccal_f$, est donnée par l'ensemble de points
		\[ \Ccal_f = \Bigset{ \bigpar{x ; f(x)} \text{ où $x$ parcourt } \D }. \]
	On lit « la courbe représentative de $f$ est l'ensemble des points $\bigl(x ; f(x)\bigr)$ du plan où $x$ parcourt le domaine de $f$ ».
}{dfn:Cf}	



\cor{propriété fondamentale}{
	La \emphindex{propriété fondamentale} suivante est valable pour tout $x, y\in\R$.
		\begin{align*}
			(x;y) \in \Ccal_f && \iff && y = f(x).
		\end{align*}
	Elle permet de décider si un point appartient à une courbe ou non.
	%\emph{Note : $f(x) = \dfrac1{200}(x+3,5)(x+2,5)(x-3,5)(x-4,5)(x-5,5) + 2$, polynôme de degré $5$ dans l'exemple ci-dessus.}
}{cor:prop-fond}



\subsection*{Résoudre graphiquement $f(x) = g(x)$}

On cherche à résoudre graphiquement une équation du type $f(x)=g(x)$, c'est-à-dire trouver l'ensemble des $x\in\R$ pour lesquels l'équation est vérifiée.

Soient $f, g$ sur $\D = ]{-}3,4 ; 2,3[$ deux fonctions réelles dont les courbes représentatives sont données ci-dessous.

\begin{figure}[h]
	\begin{center}
	 \includegraphics[page=20, scale=1.25]{2-figures.pdf}
	\end{center}
	\caption{Résolution graphique de l'équation $f(x) = g(x)$}
\end{figure}
	
Imaginons un nombre réel $x\in\R$ parcourant $]{-}3,4 ; 2,3[$ en partant de la gauche et allant vers la droite.
Pour chaque $x$, le point de $\Ccal_f$ d'abscisse $x$ est d'ordonnée $f(x)$, par définition.
Idem pour le point de $\Ccal_g$ d'abscisse $x$ : il est d'ordonnée $g(x)$.

Ainsi quand $x$ traverse le domaine, ces deux points sont confondus si et seulement si ces deux points sont de même ordonnée, c'est-à-dire si et seulement si $f(x) = g(x)$.		
	
Les solutions $x\in]{-}3,4 ; 2,3[$ de l'équation $f(x) = g(x)$ sont donc données approximativement par l'ensemble
	\[ \bigset{ -0,8 ; 1,2 }. \]

\thm{}{
	Soient $f,g$ deux fonctions réelles définies sur un même domaine $\D$.
	Les solutions dans $\D$ de l'équation $f(x) = g(x)$ sont les abscisses des points d'intersection de $\Ccal_f$ et $\Ccal_g$.
}{}

\subsection*{Résoudre graphiquement $f(x) \leq g(x)$}

On cherche à résoudre graphiquement une équation du type $f(x)=g(x)$, c'est-à-dire trouver l'ensemble des $x\in\R$ pour lesquels l'inéquation est vérifiée.

Soient $f, g$ sur $\D = ]{-}3,4 ; 2,3[$ deux fonctions réelles dont les courbes représentatives sont données ci-dessous.

\begin{figure}[h]
	\begin{center}
	 \includegraphics[page=21, scale=1.25]{2-figures.pdf}
	\end{center}
	\caption{Résolution graphique de l'équation $f(x) \leq g(x)$}
\end{figure}

Pour chaque $x$ parcourant $]{-}3,4 ; 2,3[$, la hauteur (l'ordonnée) du point de $\Ccal_f$ d'abscisse $x$ est donnée par $f(x)$, et celle du point de $\Ccal_g$ d'abscisse $x$ est donnée par $g(x)$.
Ainsi, le premier point est en dessous du second dès que $f(x) \leq g(x)$.

L'ensemble des solutions de l'inéquation $f(x) \leq g(x)$ est donc donné approximativement par
	\[ [{-}0,8 ; 1,2 ]. \]

\thm{}{
	Soient $f, g$ deux fonctions réelles définies sur un même domaine $\D$.
	Les solutions dans $\D$ de l'équation $f(x) \leq g(x)$ sont les abscisses des points de $\Ccal_f$ situés en dessous de $\Ccal_g$.
}{}

\notations{
	Soient deux intervalles $I, J \subseteq \R$. On note $I \cup J$ l'\emphindex{union} des ensembles $I$ et $J$ : ce sont les nombres réels qui appartiennent à $I$ \underline{ou} à $J$.
	\begin{center}
	 \includegraphics[page=22, scale=1.25]{2-figures.pdf}
	\end{center}
}

%
%\newpage
%\addcontentsline{toc}{section}{Exercices}
%\fancyhead[C]{\textbf{Chapitre \thechapter~— Exercices}}




%\newpage
%\addcontentsline{toc}{section}{Solutions}
%\fancyhead[C]{\textbf{Chapitre \thechapter~— Solutions}}
%
%\shipoutAnswer



